\begin{enumerate}

\item El área encerrada por la línea media es $A_m = a b$.

\item La longitud de la línea media es $s_m = 2(a + b)$.

\item La constante de torsión es $4 A_m^2 t/s_m$ y la tensión es $M_x/(2 A_m t)$.

\end{enumerate}

\noindent\textbf{Operaciones y resultados intermedios.}

\begin{enumerate}

\item Área encerrada.

\[\fitmath{y_{1}=100\cdot 60\approx 6000\,\mathrm{mm^{2}}}\]

\item Tensión tangencial.

\[\fitmath{y_{2}=\frac{1800000}{2\cdot 100\cdot 60\cdot 2}\approx 75\,\mathrm{MPa}}\]

\item Constante de torsión.

\[\fitmath{y_{3}=\frac{4\cdot {\left(100\cdot 60\right)}^{2}\cdot 2}{2\cdot \left(100+60\right)}\approx 900000\,\mathrm{mm^{4}}}\]

\item Giro entre extremos.

\[\fitmath{y_{4}=\frac{1800000\cdot 1000}{\frac{26950\cdot 4\cdot {\left(100\cdot 60\right)}^{2}\cdot 2}{2\cdot \left(100+60\right)}}\approx 0{,}0742115\,\mathrm{rad}}\]

\end{enumerate}
